\documentclass{article}
\usepackage{longtable}
\usepackage{booktabs}
\usepackage{array}

\title{\textbf{Relatório Tableaux Semântico}}
\author{Erick Schultz \\ Weverton Oliveira da Silva}

\begin{document}

\maketitle

\section{Utilização}

\section*{Execução}

No diretório raiz do projeto, execute o comando:

\begin{verbatim}
python -m src.main
\end{verbatim}

\section*{Sintaxe}

As fórmulas devem ser escritas utilizando letras como átomos
proposicionais, por exemplo, \texttt{p}, \texttt{q} e
\texttt{r}.

Os operadores lógicos aceitos são:

\begin{itemize}
    \item Negação: \texttt{\~}, \texttt{¬}
    \item Conjunção: \texttt{\&}, \texttt{$\land$}
    \item Disjunção: \texttt{|}, \texttt{$\lor$}
    \item Implicação: \texttt{->}, \texttt{=>}, \texttt{→}, \texttt{>>}
    \item Bicondicional: \texttt{<=>}, \texttt{<->}, \texttt{$\leftrightarrow$}
    \item Parênteses: \texttt{(}, \texttt{)}
\end{itemize}

$\\\\$

\section{Sequentes utilizados}

Foram utilizados 18 sequentes para testar o funcionamento do
verificador, contemplando casos válidos e inválidos.

\begin{center}
\small
\begin{longtable}{c c c}
\toprule
\textbf{Nº} & \textbf{Sequente} & \textbf{Resultado} \\
\midrule
\endfirsthead

\toprule
\textbf{Nº} & \textbf{Sequente} & \textbf{Resultado} \\
\midrule
\endhead

1  & $p,\, p \to q \vdash q$ & Válido \\
2  & $p \to q,\, \neg q \vdash \neg p$ & Válido \\
3  & $p \lor q,\, \neg p \vdash q$ & Válido \\
4  & $p \to q,\, q \to r \vdash p \to r$ & Válido \\
5  & $\neg q \to \neg p \vdash p \to q$ & Válido \\
6  & $\neg p \to \neg q \vdash p \to q$ & Inválido \\
7  & $p \to q \vdash p \to (q \lor r)$ & Válido \\
8  & $p \to q \vdash p \to (q \land r)$ & Inválido \\
9  & $\neg(p \land q) \vdash \neg p \land \neg q$ & Inválido \\
10 & $\neg(p \lor q) \vdash \neg p \lor \neg q$ & Válido \\
11 & $p \lor q,\, \neg q \vdash p$ & Válido \\
12 & $p \to q,\, q \vdash p$ & Inválido \\
13 & $p,\, (p \land q) \to r \vdash r$ & Inválido \\
14 & $p \lor q,\, p \to r,\, q \to (r \lor s) \vdash r$ & Inválido \\
15 & $p \lor q,\, p \to r,\, q \to r \vdash r$ & Válido \\
16 & $p \vdash q$ & Inválido \\
17 & $p \to q,\, \neg p \vdash \neg q$ & Inválido \\
18 & $p \lor q,\, p \vdash \neg q$ & Inválido \\

\bottomrule
\end{longtable}
\end{center}


\section{Estruturas, estratégias e dificuldades}

\section*{Estruturas de dados}

Para representar as fórmulas e os sequentes, aproveitamos a hierarquia de classes já presente em \texttt{src/code\_base\_EP1.py} (\texttt{Proposition}, \texttt{Variable}, \texttt{Constant}, \texttt{Not}, \texttt{And}, \texttt{Or}, \texttt{Implies}, \texttt{ImpliedBy}, \texttt{Iff} e \texttt{ArgumentForm}). Sobre ela, construímos o módulo \texttt{src/tableaux.py} com as seguintes estruturas principais:

\begin{itemize}
    \item \texttt{FormulaMarcada}: associa uma fórmula à sua marca booleana (\texttt{T} ou \texttt{F}), além de armazenar sua chave sintática imutável (\texttt{chave}) e o número de nós da subárvore (\texttt{tamanho}). Utilizamos \texttt{\_\_slots\_\_} para reduzir o consumo de memória durante a expansão.
    \item \texttt{ramo} e \texttt{presentes}: em vez de manter toda a árvore do tableau em memória, armazenamos apenas o ramo ativo em uma lista linear (\texttt{ramo}) acompanhada de um conjunto (\texttt{presentes}) com os identificadores \texttt{(marca, chave)} das fórmulas já inseridas. Isso permite verificar contradições (presença do par conjugado) e evitar inserções duplicadas em tempo $O(1)$.
    \item \texttt{lineares} e \texttt{betas}: separamos os índices das fórmulas ainda não expandidas no ramo em duas estruturas: uma fila (\texttt{deque}) para regras sem bifurcação e uma lista para regras que bifurcam.
    \item \texttt{pilha\_ramos}: pilha usada no \emph{backtracking}, onde cada entrada guarda a fórmula do ramo alternativo ($\beta_2$), o tamanho anterior do ramo e uma cópia dos índices de \texttt{betas} pendentes.
    \item \texttt{ResultadoTableau}: estrutura de retorno que reúne a validade do sequente, o dicionário com o contra-exemplo (quando inválido) e a lista de pares \texttt{(regra, fórmula)} aplicados durante a prova.
\end{itemize}

\section*{Estratégias utilizadas}

O provador segue o algoritmo geral de Tableaux Analíticos com três escolhas principais de estratégia:

\begin{itemize}
    \item \textbf{Busca em profundidade (DFS):} o algoritmo expande sempre o ramo atual até que ele feche ou sature. Em uma bifurcação, o segundo ramo ($\beta_2$) é empilhado em \texttt{pilha\_ramos} e a busca continua em $\beta_1$. Quando o ramo fecha, basta descartar as posições adicionadas após a última bifurcação e retomar o topo da pilha. Essa abordagem usa menos memória que a busca em largura e encontra contra-exemplos mais rapidamente quando o sequente é inválido.
    \item \textbf{Regra $\alpha$-primeiro:} todas as expansões lineares ($\alpha\land$, $\alpha\lor$, $\alpha\to$, $\alpha\neg$ e também $\beta\neg$, que não bifurca) são aplicadas antes de qualquer regra $\beta$ bifurcável. Como expansões lineares não criam novos ramos e podem fechar o ramo atual de imediato, aplicá-las primeiro reduz bastante o tamanho final do tableau.
    \item \textbf{Heurística de seleção $\beta$:} quando restam apenas fórmulas bifurcáveis, a função \texttt{\_escolher\_beta} dá preferência à fórmula cujas subfórmulas ($\beta_1$ ou $\beta_2$) fecham o ramo imediatamente com alguma fórmula já presente em \texttt{presentes}. Em caso de empate, seleciona a fórmula de menor tamanho sintático. Se apenas $\beta_2$ fechar de imediato, invertemos a ordem para explorar esse ramo primeiro.
\end{itemize}

Além disso, implementamos um parser descendente recursivo em \texttt{src/parser.py} para converter sequentes em texto diretamente para \texttt{ArgumentForm}, respeitando a precedência dos conectivos ($\neg > \land > \lor > \to > \leftrightarrow$) e a associatividade à direita na implicação.

\section*{Dificuldades encontradas}

Durante o desenvolvimento, os principais desafios foram:

\begin{itemize}
    \item \textbf{Comparação de fórmulas no código-base:} o método \texttt{\_\_eq\_\_} de \texttt{Proposition} em \texttt{code\_base\_EP1.py} testa equivalência lógica gerando toda a tabela-verdade das fórmulas. Usar essa comparação para checar contradições ou duplicatas tornaria o provador exponencial mesmo em casos simples. Resolvemos isso criando \texttt{chave\_formula}, que gera uma representação sintática em tupla para comparação direta em conjunto.
    \item \textbf{Operadores n-ários:} os conectivos \texttt{\&} e \texttt{|} do código-base achatam expressões encadeadas em um único nó com mais de dois filhos (por exemplo, \texttt{P \& Q \& R} gera um \texttt{And} com três filhos), além de aceitarem 0 ou 1 filho. Foi necessário tratar esses casos na função \texttt{expandir\_formula}, separando o primeiro filho e reagrupando os restantes.
    \item \textbf{Tratamento de $\beta\neg$ e constantes:} apesar de $F\,\neg C$ ser classificada como regra $\beta$ na tabela do enunciado, ela não gera bifurcação. Colocá-la na fila linear evitou duplicação de ramos e de regras registradas. Também foi preciso tratar explicitamente as constantes \texttt{True} e \texttt{False}, fechando o ramo sempre que surgem \texttt{T False} ou \texttt{F True}.
    
\end{itemize}
\end{document}
